\chapter{Meeting a Japanese Teacher Online}

% ================= PERCAKAPAN 1-6 =================
\begin{convobox}{1}
    \noindent
    \jpkana{あっ、}{A,,} 
    \jpkana{こんにちは。}{konnichiwa.} 

    \vspace{1.5em}
    \noindent{\Large \color{articolor} \textbf{Arti:} Ah, halo (selamat siang).}
\end{convobox}

\begin{convobox}{2}
    \noindent
    \jpword{わたし}{私}{Watashi} \jpkana{の}{no} 
    \jpword{こえ}{声}{koe}\jpkana{、}{,} 
    \jpword{き}{聞}{ki}\jpkana{こえますか？}{koemasu ka?}

    \vspace{1.5em}
    \noindent{\Large \color{articolor} \textbf{Arti:} Apakah suaraku terdengar?}
\end{convobox}

\begin{convobox}{3}
    \noindent
    \jpkana{はい、}{Hai,} 
    \jpword{き}{聞}{ki}\jpkana{こえます。}{koemasu.}

    \vspace{1.5em}
    \noindent{\Large \color{articolor} \textbf{Arti:} Ya, terdengar.}
\end{convobox}

\begin{convobox}{4}
    \noindent
    \jpkana{はじめまして、}{Hajimemashite,} 
    \jpkana{Tanaka}{Tanaka} 
    \jpkana{と}{to} 
    \jpword{もう}{申}{mou}\jpkana{します。}{shimasu.}

    \vspace{1.5em}
    \noindent{\Large \color{articolor} \textbf{Arti:} Salam kenal, namaku Tanaka.}
\end{convobox}

\begin{convobox}{5}
    \noindent
    \jpkana{えっと、}{Eetto,} 
    \jpkana{お}{o}\jpword{なまえ}{名前}{namae} \jpkana{は}{wa} 
    \jpword{なん}{何}{nan} \jpkana{ですか？}{desu ka?}

    \vspace{1.5em}
    \noindent{\Large \color{articolor} \textbf{Arti:} Emm, nama kamu siapa?}
\end{convobox}

\begin{convobox}{6}
    \noindent
    \jpkana{クロワッサン}{Kurowassan} 
    \jpkana{と}{to} 
    \jpword{もう}{申}{mou}\jpkana{します。}{shimasu.} 
    \jpkana{よろしくお}{Yoroshiku o}\jpword{ねが}{願}{nega}\jpkana{いします。}{ishimasu.}

    \vspace{1.5em}
    \noindent{\Large \color{articolor} \textbf{Arti:} Namaku Croissant. Mohon bantuannya (salam kenal).}
\end{convobox}

% --- LESSON BUFFER 1 ---
\begin{lessonbox}{Catatan Belajar: Halo, Kenalan Yuk!}
    Halo adik-adik yang pintar! Kalau kita baru pertama kali bertemu seseorang, kita mengucapkan \textbf{はじめまして (Hajimemashite)} yang artinya "Salam kenal". \\
    
    Nah, untuk menyebutkan nama, biasanya kita pakai "Namaku... desu". Tapi, guru Tanaka menggunakan kata \textbf{と申します (to moushimasu)}. Wah, apa itu? Itu adalah cara yang saaaangat sopan untuk menyebut nama kita! Jadi, kalau mau terdengar sangat keren dan sopan, kamu bisa bilang: \textit{"(Namamu) to moushimasu!"}. Jangan lupa akhiri perkenalanmu dengan \textit{Yoroshiku onegaishimasu} ya!
\end{lessonbox}

% ================= PERCAKAPAN 7-13 =================
\begin{convobox}{7}
    \noindent
    \jpkana{よろしくお}{Yoroshiku o}\jpword{ねが}{願}{nega}\jpkana{いします。}{ishimasu.} 

    \vspace{1.5em}
    \noindent{\Large \color{articolor} \textbf{Arti:} Salam kenal juga.}
\end{convobox}

\begin{convobox}{8}
    \noindent
    \jpkana{オンライン}{Onrain} \jpkana{の}{no} 
    \jpword{にほんご}{日本語}{Nihongo} 
    \jpkana{レッスン}{ressun} \jpkana{は、}{wa,} 
    \jpword{はじ}{初}{haji}\jpkana{めてですか？}{mete desu ka?}

    \vspace{1.5em}
    \noindent{\Large \color{articolor} \textbf{Arti:} Apakah ini pertama kalinya ikut pelajaran bahasa Jepang online?}
\end{convobox}

\begin{convobox}{9}
    \noindent
    \jpkana{はい、}{Hai,} 
    \jpword{はじ}{初}{haji}\jpkana{めてです。}{mete desu.}

    \vspace{1.5em}
    \noindent{\Large \color{articolor} \textbf{Arti:} Ya, ini pertama kalinya.}
\end{convobox}

\begin{convobox}{10}
    \noindent
    \jpkana{そうですか。}{Sou desu ka.} 
    \jpkana{お}{o}\jpword{あ}{会}{a}\jpkana{いできて}{i dekite} 
    \jpkana{うれしいです。}{ureshii desu.}

    \vspace{1.5em}
    \noindent{\Large \color{articolor} \textbf{Arti:} Oh begitu. Aku senang bisa bertemu denganmu.}
\end{convobox}

\begin{convobox}{11}
    \noindent
    \jpkana{いつから}{Itsu kara} 
    \jpword{にほんご}{日本語}{Nihongo} \jpkana{の}{no} 
    \jpword{べんきょう}{勉強}{benkyou} \jpkana{を}{o} 
    \jpkana{していますか？}{shite imasu ka?}

    \vspace{1.5em}
    \noindent{\Large \color{articolor} \textbf{Arti:} Sejak kapan kamu belajar bahasa Jepang?}
\end{convobox}

\begin{convobox}{12}
    \noindent
    \jpword{にねんまえ}{2年前}{Ni-nen mae} \jpkana{から}{kara} 
    \jpword{べんきょう}{勉強}{benkyou} \jpkana{しています。}{shite imasu.}

    \vspace{1.5em}
    \noindent{\Large \color{articolor} \textbf{Arti:} Aku belajar sejak 2 tahun yang lalu.}
\end{convobox}

\begin{convobox}{13}
    \noindent
    \jpkana{どうして}{Doushite} 
    \jpword{にほんご}{日本語}{Nihongo} \jpkana{の}{no} 
    \jpword{べんきょう}{勉強}{benkyou} \jpkana{を}{o} 
    \jpword{はじ}{始}{haji}\jpkana{めたんですか？}{meta n desu ka?}

    \vspace{1.5em}
    \noindent{\Large \color{articolor} \textbf{Arti:} Kenapa kamu mulai belajar bahasa Jepang?}
\end{convobox}

% --- LESSON BUFFER 2 ---
\begin{lessonbox}{Catatan Belajar: Bertanya "Kapan?" dan "Kenapa?"}
    Saat kita penasaran, kita bisa bertanya loh! Di percakapan ini, Guru Tanaka bertanya:
    \begin{itemize}
        \item \textbf{いつ (Itsu)} = Kapan?
        \item \textbf{どうして (Doushite)} = Kenapa?
    \end{itemize}
    Coba bayangkan bahasa Jepang itu seperti main tebak-tebakan. Kalau gurumu tanya \textit{"Doushite?"} (Kenapa), nanti kamu harus jawab dengan alasanmu! Yuk kita lihat apa alasannya di kotak selanjutnya!
\end{lessonbox}

% ================= PERCAKAPAN 14-20 =================
\begin{convobox}{14}
    \noindent
    \jpword{にほん}{日本}{Nihon} \jpkana{の}{no} 
    \jpword{まんが}{漫画}{manga} \jpkana{が}{ga} 
    \jpword{す}{好}{su}\jpkana{きだからです。}{ki da kara desu.}

    \vspace{1.5em}
    \noindent{\Large \color{articolor} \textbf{Arti:} Karena aku suka komik Jepang.}
\end{convobox}

\begin{convobox}{15}
    \noindent
    \jpword{にほん}{日本}{Nihon} \jpkana{に}{ni} 
    \jpword{き}{来}{ki}\jpkana{たことは}{ta koto wa} 
    \jpkana{ありますか？}{arimasu ka?}

    \vspace{1.5em}
    \noindent{\Large \color{articolor} \textbf{Arti:} Apakah kamu pernah datang ke Jepang?}
\end{convobox}

\begin{convobox}{16}
    \noindent
    \jpkana{はい、}{Hai,} 
    \jpword{いっかい}{1回}{ikkai} \jpkana{あります。}{arimasu.} 
    \jpword{とうきょう}{東京}{Toukyou} \jpkana{と}{to} 
    \jpword{おおさか}{大阪}{Oosaka} \jpkana{と}{to} 
    \jpword{なら}{奈良}{Nara} \jpkana{に}{ni} 
    \jpword{い}{行}{i}\jpkana{きました。}{kimashita.}

    \vspace{1.5em}
    \noindent{\Large \color{articolor} \textbf{Arti:} Ya, pernah 1 kali. Aku pergi ke Tokyo, Osaka, dan Nara.}
\end{convobox}

\begin{convobox}{17}
    \noindent
    \jpword{ふだん}{普段}{Fudan} 
    \jpkana{どうやって}{dou yatte} 
    \jpword{にほんご}{日本語}{Nihongo} \jpkana{を}{o} 
    \jpword{べんきょう}{勉強}{benkyou} \jpkana{していますか？}{shite imasu ka?}

    \vspace{1.5em}
    \noindent{\Large \color{articolor} \textbf{Arti:} Biasanya bagaimana cara kamu belajar bahasa Jepang?}
\end{convobox}

\begin{convobox}{18}
    \noindent
    \jpword{にほんご}{日本語}{Nihongo} \jpkana{の}{no} 
    \jpword{まんが}{漫画}{manga} \jpkana{を}{o} 
    \jpword{よ}{読}{yo}\jpkana{んで}{nde} 
    \jpword{べんきょう}{勉強}{benkyou} \jpkana{しています。}{shite imasu.}

    \vspace{1.5em}
    \noindent{\Large \color{articolor} \textbf{Arti:} Aku belajar dengan membaca komik bahasa Jepang.}
\end{convobox}

\begin{convobox}{19}
    \noindent
    \jpkana{それから、}{Sore kara,} 
    \jpkana{ときどき}{tokidoki} 
    \jpkana{YouTube}{YouTube} \jpkana{を}{o} 
    \jpword{み}{観}{mi}\jpkana{て}{te}

    \vspace{1.5em}
    \noindent{\Large \color{articolor} \textbf{Arti:} Lalu, kadang-kadang aku menonton YouTube}
\end{convobox}

\begin{convobox}{20}
    \noindent
    \jpword{かいわ}{会話}{kaiwa} \jpkana{の}{no} 
    \jpword{れんしゅう}{練習}{renshuu} \jpkana{を}{o} 
    \jpkana{しています。}{shite imasu.}

    \vspace{1.5em}
    \noindent{\Large \color{articolor} \textbf{Arti:} dan latihan percakapan.}
\end{convobox}

% --- LESSON BUFFER 3 ---
\begin{lessonbox}{Catatan Belajar: Menjawab "Karena..."}
    Nah, ini dia jawaban dari \textit{"Doushite?"} (Kenapa?) tadi! Kalau kamu mau bilang "Karena...", kamu tinggal menambahkan \textbf{〜から (kara)} di akhir kalimatmu. \\
    
    Contohnya: \textit{Manga ga suki da \textbf{kara} desu!} (Karena suka komik!). \\
    Gampang banget kan? Huruf-huruf di atas mungkin terlihat banyak, tapi pelan-pelan saja bacanya. Romaji di bawahnya akan menuntunmu!
\end{lessonbox}

% ================= PERCAKAPAN 21-27 =================
\begin{convobox}{21}
    \noindent
    \jpkana{レッスン}{Ressun} \jpkana{で}{de} 
    \jpkana{どんなことを}{donna koto o} 
    \jpkana{したいですか？}{shitai desu ka?}

    \vspace{1.5em}
    \noindent{\Large \color{articolor} \textbf{Arti:} Di pelajaran nanti, apa yang ingin kamu lakukan?}
\end{convobox}

\begin{convobox}{22}
    \noindent
    \jpword{かいわ}{会話}{Kaiwa} \jpkana{の}{no} 
    \jpword{れんしゅう}{練習}{renshuu} \jpkana{を}{o} 
    \jpkana{したいですか？}{shitai desu ka?}

    \vspace{1.5em}
    \noindent{\Large \color{articolor} \textbf{Arti:} Apakah ingin berlatih percakapan?}
\end{convobox}

\begin{convobox}{23}
    \noindent
    \jpkana{それとも、}{Sore tomo,} 
    \jpword{ぶんぽう}{文法}{bunpou} \jpkana{を}{o} 
    \jpword{べんきょう}{勉強}{benkyou} \jpkana{したいですか？}{shitai desu ka?}

    \vspace{1.5em}
    \noindent{\Large \color{articolor} \textbf{Arti:} Atau, apakah ingin belajar tata bahasa?}
\end{convobox}

\begin{convobox}{24}
    \noindent
    \jpword{かいわ}{会話}{Kaiwa} \jpkana{の}{no} 
    \jpword{れんしゅう}{練習}{renshuu} \jpkana{を}{o} 
    \jpkana{たくさん}{takusan} 
    \jpkana{したいです。}{shitai desu.}

    \vspace{1.5em}
    \noindent{\Large \color{articolor} \textbf{Arti:} Aku ingin banyak berlatih percakapan.}
\end{convobox}

\begin{convobox}{25}
    \noindent
    \jpkana{あと、}{Ato,} 
    \jpword{しゅくだい}{宿題}{shukudai} \jpkana{も}{mo} 
    \jpkana{たくさん}{takusan} 
    \jpword{だ}{出}{da}\jpkana{してほしいです。}{shite hoshii desu.}

    \vspace{1.5em}
    \noindent{\Large \color{articolor} \textbf{Arti:} Terus, aku ingin diberikan banyak PR juga.}
\end{convobox}

\begin{convobox}{26}
    \noindent
    \jpword{きょうかしょ}{教科書}{Kyoukasho} \jpkana{は、}{wa,} 
    \jpkana{何か}{nani ka} 
    \jpword{も}{持}{mo}\jpkana{っていますか？}{tte imasu ka?}

    \vspace{1.5em}
    \noindent{\Large \color{articolor} \textbf{Arti:} Apakah kamu punya buku pelajaran?}
\end{convobox}

\begin{convobox}{27}
    \noindent
    \jpkana{いえ、}{Ie,} 
    \jpword{なに}{何}{nani}\jpkana{も}{mo} 
    \jpword{も}{持}{mo}\jpkana{っていません。}{tte imasen.} 
    \jpkana{おすすめの}{Osusume no} 
    \jpword{きょうかしょ}{教科書}{kyoukasho} \jpkana{は}{wa} 
    \jpkana{ありますか？}{arimasu ka?}

    \vspace{1.5em}
    \noindent{\Large \color{articolor} \textbf{Arti:} Tidak, aku tidak punya apa-apa. Apakah ada buku pelajaran yang direkomendasikan?}
\end{convobox}

% --- LESSON BUFFER 4 ---
\begin{lessonbox}{Catatan Belajar: "Aku Ingin..."}
    Tahukah kamu cara bilang "Aku ingin..." dalam bahasa Jepang? Gampang banget! Kita cukup menambahkan kata \textbf{〜たい (〜tai)} di belakang kegiatannya!\\
    
    Contoh:
    \begin{itemize}
        \item Belajar = \textit{Benkyou shimasu} $\rightarrow$ Ingin belajar = \textbf{Benkyou shitai desu}.
    \end{itemize}
    Wow, teman kita ini rajin sekali ya, sampai minta banyak \textbf{宿題 (Shukudai)} alias PR! Adik-adik suka PR juga nggak nih?
\end{lessonbox}

% ================= PERCAKAPAN 28-35 =================
\begin{convobox}{28}
    \noindent
    \jpkana{そろそろ}{Sorosoro} 
    \jpword{じかん}{時間}{jikan} \jpkana{なので、}{na node,} 
    \jpword{きょう}{今日}{kyou} \jpkana{は}{wa} 
    \jpkana{ここまでにしましょう。}{koko made ni shimashou.}

    \vspace{1.5em}
    \noindent{\Large \color{articolor} \textbf{Arti:} Karena sudah hampir waktunya, hari ini kita akhiri sampai di sini ya.}
\end{convobox}

\begin{convobox}{29}
    \noindent
    \jpword{さいご}{最後}{Saigo} \jpkana{に、}{ni,} 
    \jpkana{何か}{nani ka} 
    \jpword{しつもん}{質問}{shitsumon} \jpkana{は}{wa} 
    \jpkana{ありますか？}{arimasu ka?}

    \vspace{1.5em}
    \noindent{\Large \color{articolor} \textbf{Arti:} Terakhir, apakah ada pertanyaan?}
\end{convobox}

\begin{convobox}{30}
    \noindent
    \jpkana{いえ、}{Ie,} 
    \jpword{だいじょうぶ}{大丈夫}{daijoubu} \jpkana{です。}{desu.}

    \vspace{1.5em}
    \noindent{\Large \color{articolor} \textbf{Arti:} Tidak, sudah cukup (tidak ada).}
\end{convobox}

\begin{convobox}{31}
    \noindent
    \jpword{きょう}{今日}{Kyou} \jpkana{の}{no} 
    \jpkana{レッスン}{ressun} \jpkana{は、}{wa,} 
    \jpkana{どうでしたか？}{dou deshita ka?} 
    \jpword{むずか}{難}{muzuka}\jpkana{しかったですか？}{shikatta desu ka?}

    \vspace{1.5em}
    \noindent{\Large \color{articolor} \textbf{Arti:} Pelajaran hari ini bagaimana? Apakah sulit?}
\end{convobox}

\begin{convobox}{32}
    \noindent
    \jpword{にほんご}{日本語}{Nihongo} \jpkana{で}{de} 
    \jpword{はな}{話}{hana}\jpkana{すのは}{su no wa} 
    \jpword{むずか}{難}{muzuka}\jpkana{しかったですが、}{shikatta desu ga,} 
    \jpword{たの}{楽}{tano}\jpkana{しかったです。}{shikatta desu.}

    \vspace{1.5em}
    \noindent{\Large \color{articolor} \textbf{Arti:} Berbicara dengan bahasa Jepang itu sulit, tapi menyenangkan.}
\end{convobox}

\begin{convobox}{33}
    \noindent
    \jpword{じかい}{次回}{Jikai} \jpkana{の}{no} 
    \jpkana{レッスン}{ressun} \jpkana{は、}{wa,} 
    \jpword{なんようび}{何曜日}{nan-youbi} \jpkana{の}{no} 
    \jpword{なんじ}{何時}{nan-ji} \jpkana{ごろが}{goro ga} 
    \jpkana{いいですか？}{ii desu ka?}

    \vspace{1.5em}
    \noindent{\Large \color{articolor} \textbf{Arti:} Pelajaran berikutnya, enaknya hari apa dan kira-kira jam berapa?}
\end{convobox}

\begin{convobox}{34}
    \noindent
    \jpword{へいじつ}{平日}{Heijitsu} \jpkana{は}{wa} 
    \jpword{しごと}{仕事}{shigoto} \jpkana{なので、}{na node,} 
    \jpword{どようび}{土曜日}{doyoubi} \jpkana{の}{no} 
    \jpword{あさ}{朝}{asa} \jpword{じゅうじ}{10時}{juu-ji} \jpkana{でもいいですか？}{demo ii desu ka?}

    \vspace{1.5em}
    \noindent{\Large \color{articolor} \textbf{Arti:} Karena hari biasa aku bekerja, apakah boleh hari Sabtu pagi jam 10?}
\end{convobox}

\begin{convobox}{35}
    \noindent
    \jpkana{わかりました。}{Wakarimashita.} 
    \jpkana{では、}{Dewa,} 
    \jpkana{また}{mata} 
    \jpword{じかい}{次回}{jikai} \jpkana{の}{no} 
    \jpkana{レッスン}{ressun} \jpkana{で}{de} 
    \jpkana{お}{o}\jpword{あ}{会}{a}\jpkana{いしましょう。}{i shimashou.} 
    \jpword{しつれい}{失礼}{shitsurei} \jpkana{します。}{shimasu.}

    \vspace{1.5em}
    \noindent{\Large \color{articolor} \textbf{Arti:} Aku mengerti. Kalau begitu, sampai jumpa lagi di pelajaran berikutnya. Permisi.}
\end{convobox}

% --- SUMMARY BOX ---
\vspace{2em}
\begin{tcolorbox}[
    colback=blue!5!white,
    colframe=blue!80!black,
    boxrule=3pt,
    arc=5mm,
    title={\huge\bfseries \sffamily 🌟 Rangkuman Bab 4 🌟},
    fonttitle=\color{white},
    attach boxed title to top center={yshift=-4mm},
    boxed title style={colback=blue!80!black, rounded corners},
    left=5mm, right=5mm, top=7mm, bottom=5mm
]
\Large \textbf{Hebat sekali! Kamu sudah berani ikut kelas online!} Mari kita ingat-ingat lagi apa saja yang sudah kita pelajari:
\begin{itemize}
    \item \textbf{Menyebutkan nama dengan super sopan:} \textit{\dots to moushimasu} (Namaku adalah\dots)
    \item \textbf{Kata tanya penting:} \textit{Itsu} (Kapan) dan \textit{Doushite} (Kenapa).
    \item \textbf{Menjawab "Karena":} \textit{\dots kara desu}.
    \item \textbf{Bilang "Aku ingin\dots":} Tambahkan akhiran \textit{\dots tai desu}.
    \item \textbf{Menanyakan hari dan jam:} \textit{Nan-youbi?} (Hari apa?) dan \textit{Nan-ji?} (Jam berapa?).
\end{itemize}
Semakin lama belajar, rasanya semakin seru ya! Terus berlatih agar lidahmu semakin terbiasa mengucapkan huruf-huruf bahasa Jepang. Semangat, kamu pasti bisa!
\end{tcolorbox}

\newpage